\section{Related Work}

This framework builds upon and relates to several established approaches in philosophy of mind and consciousness studies, as well as computer science and quantum mechanics:

\subsection{Panpsychism and Russellian Monism}

Panpsychism proposes that consciousness is a fundamental feature of all matter. Even electrons and atoms have some form of basic experience. Russellian monism suggests that physics tells us what matter does (its behavior) but not what it is (its intrinsic nature), proposing that this hidden nature might be experiential. Both approaches try to solve the hard problem by putting experience at reality's foundation rather than having it mysteriously emerge from non-conscious matter.

Vibe Theory shares this starting point but goes further. While panpsychism adds experiential properties to otherwise physical entities (atoms have mass AND experience), Vibe Theory eliminates the physical substrate entirely. There are only experiential units (vibes) whose patterns create what we call matter. It solves panpsychism's ``combination problem'' (how do micro-experiences combine into human consciousness?) through hierarchical nesting where parent vibes naturally experience their children. Unlike panpsychism's often vague dynamics, Vibe Theory provides specific mechanisms: beats for time, tones for states, and links for relationships.

\subsection{Integrated Information Theory}

IIT (Tononi, 2008) proposes that consciousness corresponds to integrated information. According to IIT, any system that integrates information in a way that is greater than the sum of its parts possesses consciousness to that degree. A simple sensor would have minimal consciousness, while a human brain generates high consciousness through its richly interconnected networks. IIT attempts to solve the hard problem by providing a mathematical framework that predicts which physical systems are conscious and to what degree.

Vibe Theory differs fundamentally from IIT in its starting assumptions. While IIT begins with physical systems and asks how much consciousness they generate, Vibe Theory starts with experience as the only reality. In IIT, consciousness emerges from certain configurations of matter. In Vibe Theory, what we call matter emerges from configurations of experience. IIT measures consciousness levels in physical systems, implying some systems have zero consciousness. Vibe Theory asserts that everything is experiential. There is no non-conscious substrate from which consciousness emerges. The question isn't ``how much consciousness does this system have?'' but rather ``how is consciousness organized in this region of the experiential mesh?''

\subsection{Process Philosophy and Neutral Monism}

Whitehead's process philosophy (1929) revolutionized metaphysics by proposing that reality consists not of static substances but of ``actual occasions'' of experience. Each moment of reality is a drop of experience that feels its past, synthesizes influences, and contributes to the future. Whitehead saw the universe as composed of these experiential events rather than inert matter. Vibe Theory builds on this foundation: vibes are similar to Whitehead's occasions to some degree but with precise computational rules. Where Whitehead remained abstract about how occasions interact, Vibe Theory specifies exact mechanisms through links, beats, and update functions.

Russell's neutral monism (1927) proposed that both mind and matter emerge from a more fundamental neutral substance that is neither mental nor physical. This ``neutral stuff'' organizes differently to create mental or physical properties. Vibe Theory diverges here: rather than neutral substance, the fundamental reality is explicitly experiential. Where Russell sought a middle ground between mind and matter, Vibe Theory boldly claims experience as the only substance, with physical properties emerging from experiential patterns. This move from neutrality to experience reflects insights from meditation and direct investigation of consciousness.

\subsection{Digital Physics and Computational Universe Theories}

Digital physics (Wheeler, 1990, Fredkin, 2003, Wolfram, 2002) explores the hypothesis that physical reality can be modeled as discrete, rule-based informational processes, though proponents differ on whether computation is existentially fundamental or merely descriptive. From that perspective, the universe is considered essentially a giant computer processing information according to simple rules. Wheeler's famous ``it from bit'' suggests all physical entities emerge from yes/no binary choices. Fredkin envisions reality as a cellular automaton where space, time, and matter arise from cells updating their states. Wolfram's recent hypergraph work models the universe as nodes and edges evolving through graph rewriting rules, with space itself emerging from the network structure.

Vibe Theory shares digital physics' discrete, rule-based foundation but makes a crucial substitution: instead of abstract information or computation, the fundamental units are experiential. Where digital physics sees bits flipping between 0 and 1 with no inherent meaning, Vibe Theory sees vibes oscillating between pleasure and pain, states with intrinsic experiential quality. Digital physics asks what fundamental processes underlie physical law and often frames the answer in terms of information, computation, or rule-based evolution, without consensus on whether these are systematically fundamental or merely descriptive. Vibe Theory asks the same question and proposes experience as the fundamental primitive. The universe isn't just processing data. It's processing feelings, with consciousness built into the very fabric rather than emerging mysteriously from meaningless computation.

\subsection{Computational Architectures}

Two computational paradigms have particularly influenced the conceptualization of vibe mesh dynamics: associative computing and hyperbolic cellular automata.

Potter's associative computing paradigm (1992) provides a model that could potentially inform how massively parallel processing might work in the hypothesized vibe mesh. In associative computing, data is accessed not by address but by content, with all processing elements simultaneously comparing their contents to a broadcast pattern. Within the proposed framework, it is speculated that vibes might update in an analogous manner: each vibe simultaneously processing signals from all its neighbors, computing its next state in parallel with every other vibe. The associative model's emphasis on pattern matching and content-addressable memory suggests possible mechanisms for how vibes could recognize and respond to experiential patterns in their local neighborhood. It is conceivable that, just as associative processors can perform complex searches and transformations through simple parallel operations, the vibe mesh might generate consciousness through massively parallel experiential processing.

Margenstern's work on cellular automata in hyperbolic spaces (2013) suggests intriguing possibilities for vibe mesh geometry within the speculative framework. While traditional cellular automata use Euclidean grids, hyperbolic space allows exponential growth in neighbors as distance increases. This property could potentially explain how vibe meshes might support both local clustering and long-range connections. In hyperbolic geometry, a vibe could hypothetically have a small set of immediate neighbors yet be only a few links away from vast numbers of other vibes. The implementation techniques Margenstern develops for simulating hyperbolic CA on conventional computers could possibly guide practical approaches to modeling the proposed vibe dynamics, though this remains highly speculative.

\subsection{Quantum Cellular Automata}

Quantum cellular automata (QCA) extend classical cellular automata by allowing cells to exist in superposition states, updating through unitary evolution rather than deterministic rules. Unlike classical CAs where a cell is definitively alive or dead, QCA cells can be in quantum superpositions of states. This provides a natural bridge to the vibe framework: vibes might not simply be in pleasure OR pain states, but could exist in superpositions during their update cycle, only collapsing to definite tones upon interaction. The way QCA preserves quantum coherence while enabling local interactions mirrors how vibe meshes might maintain experiential unity while allowing individual vibes to evolve. Some research in quantum cellular automata suggests that certain quantum field–like behaviors may be approximated or modeled in limited regimes or toy models using discrete quantum update rules, though a full derivation of quantum field theory remains an open problem.

't Hooft's cellular automaton interpretation of quantum mechanics (2016) provides particularly intriguing support for computational approaches to fundamental physics. While not mainstream, the work demonstrates that even quantum mechanics itself could potentially emerge from deterministic cellular automata operating at the Planck scale, and this line of thinking opened up ideas on how CAs and more general CAs could potentially be used as the basic conceptual model for the vibe mesh in the framework. In 't Hooft's model, quantum superpositions arise from our coarse-grained observations of fundamentally discrete, deterministic processes. This resonates with the speculative Vibe Theory framework in suggesting that what appears continuous and probabilistic at one scale might emerge from discrete, deterministic rules at a more fundamental level.

\subsection{Graph Grammars and Topological Dynamics}

Graph grammars provide formal rules for transforming network structures, which could provide formal tools suitable for describing how links form, dissolve, and reorganize. Just as linguistic grammars generate valid sentences from simple rules, graph grammars could generate valid vibe configurations from basic linking principles. A grammar rule might specify: ``when three vibes form a triangle of mutual pain, they spontaneously create a fourth vibe in the center experiencing pleasure.'' Such rules, applied recursively, could generate the entire cosmic mesh from minimal seeds. Causal set theory adds another layer by treating spacetime itself as a discrete network of events. In the vibe framework, these ``events'' would be vibe updates, and the causal structure would emerge from the pattern of influences flowing through links. These mathematical frameworks offer formal tools for modeling evolving relational structures, which could potentially be adapted to describe experiential networks within the vibe framework.

\subsection{Holographic Principle and AdS/CFT}

The holographic principle (Susskind, 1995, Maldacena, 1997) suggests reality in a volume can be encoded on its boundary. This might relate to how higher-dimensional vibe structures could be encoded in lower-dimensional patterns, or how the mesh's global properties emerge from local interactions. The AdS/CFT correspondence particularly suggests deep connections between seemingly different descriptions of reality.

\subsection{Quantum Error Correction and Information Preservation}

Holographic quantum error-correcting codes (Pastawski et al., 2015) show how quantum information in a bulk space can be redundantly encoded on a boundary, protecting against local errors. Within the speculative vibe framework, this could suggest how stable experiences might stay coherent despite local disruptions in the vibe mesh. Hierarchical vibe mesh networks could potentially implement a form of experiential error correction, which could allow for higher-order structures to build on top, though this remains highly conjectural.

\subsection{Neural Networks and Artificial Intelligence}

Neural networks parallel the vibe mesh: neurons sum weighted inputs like vibes integrating neighbor influences, hierarchies mirror nested vibe structures, and attention mechanisms resemble variable link strengths. Hopfield networks show how simple rules create stable patterns. The success of AI demonstrates how local computational units can generate complex intelligent behavior, though current networks lack the experiential qualities central to vibes. If we could engineer artificial vibes with genuine experiential states rather than mere information processing, we might create truly conscious AI that feels rather than merely computes.

\subsection{Relational Quantum Mechanics}

Rovelli's relational quantum mechanics (1996) eliminates observer-independent states, proposing that quantum states exist only relative to observers. This aligns with this relational view where vibes exist only through their connections. However, RQM maintains a physical substrate, whereas Vibe Theory makes the relations themselves experiential. QBism (Fuchs et al., 2014) goes further by putting agent experience at the center of quantum mechanics, approaching this experiential foundation.

\subsection{Religious and Spiritual Traditions}

Vibe Theory can be considered to touch on similar themes to many of the worlds ancient religious or spiritual traditions as well. Judaism's Kabbalistic tzimtzum parallels how unity differentiates into multiplicity. Buddhism's Indra's Net and dependent origination mirror the interconnected vibe mesh. Christianity's Holy Spirit as animating presence aligns with vibes as experiential units. Hinduism's concept of spanda (vibratory consciousness) bears a thematic resemblance to the vibe concept. Islam's tawhid (absolute unity of God) and Sufi concepts of wahdat al-wujud (unity of being) reflect the underlying oneness. Taoism's qi and yin-yang dynamics resemble vibe flows and tone polarities. Indigenous traditions worldwide speak of reality as a living web of relationships and interconnectedness. These traditions and others express, through word and practice, themes of unity, relationality, and dynamic becoming that this framework explores through abstract formal and computational metaphors.

\subsection{Reaction-Diffusion Systems and Pattern Formation}

Reaction–diffusion systems provide a useful analogy for how local interactions can generate global patterns (Turing patterns). This parallels how simple vibe update rules might create complex experiential structures from initially uniform states.

\subsection{Loop Quantum Gravity and Discrete Spacetime}

Loop quantum gravity's use of discrete relational structures offers a conceptual parallel to the vibe mesh, despite differing existential or ontological commitments. The mathematical tools from LQG might help formalize vibe dynamics.

\subsection{Monadology and Occasionalism}

Leibniz's monads can be interpreted as early philosophical proposals of simple, irreducible units organized into hierarchical structures, bearing partial resemblance to the role vibes play in this framework. Unlike monads' pre-established harmony, vibes interact directly through links. Islamic occasionalism similarly explored discrete moments forming continuous experience.

\subsection{Quantum Field Theory as Computation}

Some authors have explored viewing physical law as computational (Lloyd, 2006, Tegmark, 2014), though typically without assigning intrinsic experiential properties to computation itself. However, while they see computation generating consciousness, Vibe Theory sees consciousness/experience as the substrate performing computation. The universe computes its next experiential state, not its next physical state. This suggests quantum fields themselves might be patterns of vibe activity, with particles as localized computational processes within the experiential mesh.

\subsection{Comparison to a Few Other Theories}

The vibe framework relates to several philosophical positions. Whitehead's process philosophy shares the focus on experiential events, though vibes provide concrete computational mechanisms. Unlike neutral monism (James, Russell), which sees mental and physical as co-equal aspects, vibes make experience fundamental with physics emerging from it. Cosmopsychism struggles with relating universal to individual minds. The hierarchical vibe mesh solves this through parent-child nesting. Quantum consciousness theories (Penrose, Stapp) share interest in physics-consciousness connections but focus on neuroscience rather than the fundamental nature of reality.

This framework has been synthesizing and integrating aspects and ideas of these systems where it opens doors, while also introducing new thoughts of experiential dynamics and relational structure, though it remains philosophical speculation lacking empirical grounding.
